% ============================================================================
% ABSTRACT tiếng Anh — cuối bài, sau Tài liệu tham khảo (theo mẫu bài tiếng Việt):
%   - "ABSTRACT" 11pt; tên bài tiếng Anh chữ thường, in đậm, 11pt.
%   - Tác giả 12pt; đơn vị in nghiêng 11pt; tóm tắt 11pt.
% ============================================================================

\begin{english}
\vspace{4pt}

\begin{center}
{\fontsize{11}{13}\bfseries\selectfont ABSTRACT}
\end{center}
\vspace{-6pt}
\begin{center}
{\fontsize{11}{13}\bfseries\selectfont
Dynamic estimation of student ability based on Item Response Theory (IRT) and Longitudinal Structural Equation Modeling (LSEM)}
\end{center}
\vspace{-6pt}
\begin{center}
{\fontsize{12}{14}\selectfont
Nguyen Tien Dat\textsuperscript{1}}
\end{center}
\vspace{-6pt}
\begin{center}
{\fontsize{11}{13}\itshape\selectfont
\textsuperscript{1}Faculty of Information Technology, School of Mathematics and Information Technology,\\
Hanoi National University of Education, Hanoi, Vietnam}
\end{center}

{\fontsize{11}{13}\selectfont
\noindent Item Response Theory (IRT) systems deployed at scale, including Vietnam's OLM platform, typically assume a static student ability~\(\theta\), yet large-scale online learning produces densely longitudinal data that contradicts this assumption. We construct and validate an integrated IRT--LSEM framework for estimating dynamic ability~\(\theta_{it}\) on real OLM Mathematics data covering 24{,}213 high-school students and 122{,}957 assessment occasions across grades 10--12. The framework couples a 1PL measurement layer (EAP scoring) with four structural layers (LGCM, LCSM, multilevel growth, continuous-time DSEM) and a Kalman smoother that exploits EAP standard errors as time-varying observation noise. The dynamic model beats the static-Rasch baseline by \(\Delta\)AIC = 82--167 (LRT \(p \le 10^{-19}\)); LCSM reveals catch-up at grade 10 (\(\beta = -0.069\), CI excludes 0) shifting to plateau at grade 12; Kalman cuts mean EAP error variance by 32--41\% with no extra items.\par
\noindent\textbf{\textit{Keywords:}} Item Response Theory, longitudinal structural equation modeling, Kalman smoothing, dynamic ability estimation, educational data mining.
\par}
\end{english}
